This work investigated whether replacing the monolithic \textit{embedding} of DGI with decoupled and specialized representations could produce a base more suitable for the tasks of POI Category Classification and Next-POI Prediction. The results obtained in the three evaluated states indicate that it can. By incorporating a continuous spatial \textit{encoder}, a temporal \textit{encoder} (Time2Vec), and a hierarchical categorical \textit{encoder} (HGI), ST-MTLNet consistently outperforms the \textit{baseline} based only on DGI.

In the POI Category Classification task, the proposed variants outperform MTLNet in all 21 category-state combinations, with average gains of 20 to 24 percentage points. In the Next-POI Prediction task, the performance is also mostly favorable to ST-MTLNet, which wins in 16 of the 21 evaluated combinations, with emphasis on categories such as \textit{Food}. Taken together, these results show that decomposing the representation into spatial, temporal, and categorical components allows capturing patterns that DGI, by itself, does not model sufficiently.

The comparison between SIREN and Sphere2Vec-M also shows that there is no single universally superior spatial \textit{encoder}. SIREN presents better performance more frequently in Florida and California, while Sphere2Vec-M stands out in Texas, suggesting that the suitability of the \textit{encoder} depends on the geographic distribution of the POIs in each territory. Thus, the main gain of this work lies not only in the choice of a specific spatial architecture, but in the very adoption of decoupled representations as an alternative to the monolithic paradigm of DGI.

A relevant limitation remains in the \textit{Travel} category of the Next-POI Prediction task, in which the original MTLNet still obtains the best results in part of the scenarios. This suggests that long-distance movements and topological relationships between regions are still better captured by the graph structure used by DGI. In addition, since the three proposed components are used together, this work does not isolate the individual contribution of each \textit{encoder}, which restricts a more detailed analysis of the relative weight of each source of information.

Another limitation is related to the exclusive use of Gowalla, a dataset commonly used in the literature but composed of mobility records collected between February 2009 and October 2010. Consequently, the results should be interpreted with caution regarding their generalization to current urban mobility patterns.

As future work, we highlight the investigation of more flexible multi-task architectures, the exploration of more sophisticated fusion strategies between the \textit{encoders}, and the combination of graph-based representations and continuous coordinate encodings. These directions may broaden the observed gains and help reduce the limitations still present in categories such as \textit{Travel}.

% Como trabalhos futuros, destacam-se a realização de estudos de ablação para quantificar separadamente a contribuição dos componentes espacial, temporal e categórico, a investigação de estratégias de fusão mais sofisticadas entre os \textit{encoders} e a combinação entre representações baseadas em grafo e codificações contínuas de coordenadas. Também se pretende expandir a comparação experimental com modelos recentes do estado da arte para predição do próximo POI, como abordagens que incorporam sinais temporais explícitos. Essas direções podem ampliar os ganhos observados e ajudar a reduzir as limitações ainda presentes em categorias como \textit{Travel}.